\documentclass[10pt]{jsarticle}
\usepackage{ascmac,amsmath,amssymb,amsthm,bm,physics,arydshln,mathcomp,wrapfig}

\begin{document}
  \LARGE{単項式$0$の次数は「定義しない」、もしくは「$-\infty$」とする.}
  \bigskip

  \LARGE{なぜ$0$の次数は$0$ではないのか}

  \normalsize
  次数（degree）に関して成り立っていてほしい主張（以下）が成り立たなくなるため.

  $\deg(fg) = \deg f + \deg g$

  例えば, \ $f(x) = x^n, \ g = 0$ とすると $\deg(fg) = 0$ であるが, \\ $\deg f + \deg g = n + 0 = n$ となり上の主張が成立しない.

  \bigskip
  \bigskip

  \LARGE{では$-\infty$ ならOKなのか?}

  \normalsize
  先ほどと同様の$f, g$ に対して $\deg(0) = -\infty$ の場合, \\
  $\deg(fg) = \deg(0) = -\infty, \ \deg f + \deg g = n + (-\infty) = -\infty$ となるため主張が成立する.
  
\end{document}
